\chapter{Foundations and conventions}
\label{chap:foundations}
\section{Declared prior mathematics}
We use elementary logic, sets, maps, equivalence relations, quotients of sets,
induction and finite counting. The axiom of choice, in the form of Zorn's
lemma, is available when constructing bases or algebraic closures. A finite
set has its usual natural-number cardinality. Infinite cardinalities are
not replaced by a default natural number.

The systems $\N,\Z,\Q,\R,\C$ and elementary arithmetic are prior
foundations. We use the well-ordering of $\N$, integer division with
remainder, and the usual ordered-field properties of $\R$. Algebraic
statements about divisibility, principal ideals and polynomial factorization
are established in the relevant chapters. The elementary integer facts
needed before those chapters have the following independent proofs.

\begin{theorem}[Integer arithmetic]\label{thm:integer-arithmetic}
Integers $a,b$, not both zero, have a positive gcd $d$ expressible as
$d=xa+yb$ with $x,y\in\Z$. A positive prime dividing a product
divides one factor. Every positive integer has a unique factorization
into positive primes, up to their order.
\end{theorem}
\begin{proof}
The positive integer linear combinations of $a,b$ form a nonempty
subset of $\N$: one of $a,-a,b,-b$ is positive. Choose its least
member $d=xa+yb$. Divide $a=qd+r$ with $0\le r<d$. The remainder
is another integer combination of $a,b$, so minimality forces $r=0$.
Division of $b$ gives the same conclusion. Every common divisor of
$a,b$ divides $d$, proving the gcd and Bezout assertions. For $a=b=0$
use gcd zero. If a positive prime $p$ does not divide $a$, its gcd
with $a$ is one, since a positive divisor of $p$ is one or $p$.
A Bezout equation $xp+ya=1$, multiplied by $b$, shows that
$p\mid ab$ then implies $p\mid b$. This proves the prime-product
assertion. Strong induction factors every integer $n>1$: if it is
not prime, write it as the product of two positive integers strictly
between one and $n$, and factor each. For uniqueness, a prime in one
factorization divides a prime in the other by repeated application
of the product assertion, so they are equal. Cancel the common
nonzero factor and induct. The integer one has the empty
factorization; no positive-prime factorization is assigned to zero.
\end{proof}

The D05 linear-algebra foundation comprises vector spaces over a field,
bases, linear maps, subspaces, quotients, finite-dimensional dimension
identities, matrices, determinants, characteristic polynomials and trace.
Its fixed mathematical source is the completed D05 unified notes; its
accepted formal library is MathGenome v0.1.1. We recall the particular
facts needed below so that a proof invoking one has a precise meaning.

\begin{theorem}[Linear-algebra foundation]\label{thm:linear-foundation}
Let $K$ be a field.
\begin{enumerate}[label=(\roman*)]
\item Every linearly independent subset of a vector space extends to a
basis. Consequently every subspace has a linear complement and every
surjective linear map admits a linear section.
\item For $T:V\to W$, the induced map $V/\ker T\to\im T$ is a linear
isomorphism. In finite dimension,
$\dim V=\dim\ker T+\dim\im T$.
\item For $U\leq V$ in finite dimension,
$\dim V=\dim U+\dim(V/U)$; a proper subspace has smaller dimension.
An endomorphism is injective if and only if it is surjective.
\item For finite square matrices, $\det(AB)=\det A\det B$,
$\Tr(AB)=\Tr(BA)$, and determinants and traces are invariant under
change of basis. Trace is additive on block diagonal matrices.
\item A nonzero finite-dimensional vector space over an
algebraically closed field has the usual eigenvalue property: for every
$T\in\End_K(V)$ with $V\ne0$, some $\lambda\in K$ makes
$T-\lambda\id$ noninjective.
\end{enumerate}
\end{theorem}
These are declared prior results, not new assertions left for an abstract
algebra proof to establish. For clarity, the last item follows from the
nonconstant characteristic polynomial: choose a root $\lambda$ and use
$\det(T-\lambda I)=0$. The general basis assertion also has a short
set-theoretic construction. Order independent supersets of a given
independent set by inclusion. A union of a chain is independent because
each finite linear relation lies in one member of the chain. Zorn gives
a maximal independent set; a vector outside its span could be adjoined,
so it spans. Extending a subspace basis gives a complement and the
corresponding projection.

We distinguish an algebraically closed field in algebraic statements from
the analytic realization $\C$. Its algebraic closedness is proved internally by the algebraic
fundamental theorem of algebra in Appendix~\ref{app:algebraic-fta}.
That proof uses the abstract algebraic-closure construction, finite Galois
theory and the real intermediate value theorem; it does not assume
complex algebraic closedness or the transcendence appendix. The intermediate value theorem and Rolle's
theorem for real polynomials are declared prior real-calculus facts;
they are used to count real roots in the explicit Galois examples.
Geometric applications use the usual Euclidean coordinate distance,
equations of lines and circles, and the circle-area formula $\pi r^2$;
the constructions needed beyond those formulas are proved explicitly.
The circle-squaring application uses the
special transcendence theorem proved in Appendix~\ref{app:pi-transcendence}.
That appendix uses the elementary exponential, Euler's identity and the
real fundamental theorem of calculus, all declared prior analytic facts;
it does not take a transcendence theorem as a premise. These analytic
ingredients are kept distinct from the algebraic constructibility criterion.

\section{Notation that changes mathematical meaning}
Groups are multiplicative unless additive notation is stated. For
permutations, $(\sigma\tau)(x)=\sigma(\tau(x))$. Subgroups use the
inherited identity. Quotient-group operations require normality.
All rings in the main unqualified track are associative with a two-sided
identity. The zero ring is permitted; a field or domain is nontrivial.
Ring homomorphisms preserve the identity. A noncommutative quotient by
an ideal requires a two-sided ideal. A prime or maximal ideal is proper.

A polynomial is a formal finite coefficient sequence, not its evaluation
function. Its zero degree is $-\infty$ when a degree with a bottom value
is needed; statements involving an ordinary natural-number degree specify
the zero case separately. A factorization is unique up to units and order
until a normalization is imposed.

An extension $L/K$ includes a chosen embedding and scalar action. Several
extensions in a tower use compatible embeddings. Finite degree means
finite vector-space dimension, not a finite underlying set. A splitting
field or algebraic closure is unique up to an appropriate isomorphism,
not literal equality of unrelated types.

\section{How to read a proof dependency}
References to labels in these notes carry the exact statement and
hypotheses at that label. Use of a prior linear-algebra fact means
Theorem~\ref{thm:linear-foundation} or an explicitly identified D05
result. Standard set-theoretic constructions use the first section of
this chapter. All other algebraic prerequisites are developed in the
manuscript. A result stated later may be cited earlier only as a
forward reference in exposition; its proof must not depend back on the
result that invokes it. The accompanying dependency review checks that
such references do not form a proof cycle.
